\section{n3. Nonisolated nonnormal CE in every dimension over algebraically closed characteristic zero; explicit all-characteristic cusp envelope error \textless{}=2/m and char-zero non-FG.}
\textit{Supports: G331, n1, n2.}

New higher-dimensional nonnormal continuous-envelope sector and explicit quantitative example, using ENTIRE qualified G331 and own n1--n2. Let H be a proper numerical semigroup (a finitely generated additive subsemigroup of the nonnegative integers containing zero with finite complement) and let c be its conductor integer: every integer \textgreater{}=c belongs to H and c-1 does not. For any field k and d\textgreater{}=2, set R\_H=k[t\textasciicircum{}h (h in H),z\_1,...,z\_(d-1)] localized at the origin, and S=k[t,z\_1,...,z\_(d-1)] localized there. Both rings are excellent finite-type domains of dimension d. S is finite over R\_H since t is integral (some positive t\textasciicircum{}h lies downstairs), and their fraction fields coincide because gcd(H)=1. S is regular and integrally closed, hence is the normalization. There is just one centre over the origin. Its conductor is exactly t\textasciicircum{}c S. To verify this on polynomials, all terms of t-degree \textgreater{}=c multiply into the semigroup ring. If a polynomial has a term of degree h\textless{}c, multiply by t\textasciicircum{}(c-1-h); its coefficient of the forbidden degree c-1 remains nonzero. Distinct t-degrees cannot cancel it. Conductor commutes with localization because S/R\_H is finite. The reduced conductor quotients are BOTH k[z\_1,...,z\_(d-1)] at the origin, and their map is the identity.

The infinite-conductor boundary is nonempty: a valuation killing t and centered on the z-ring belongs to it. Thus n2 proves a full boundary homeomorphism including an actual positive-dimensional-support conductor stratum. In characteristic zero over algebraically closed k, the exact smooth source theorem of n1 gives a greatest continuous independently geometric C/D envelope for EVERY finite continuous signed obstacle on EVERY R\_H, in EVERY dimension d. Uniform convergence holds in every canonical degree and for every primary gauge. These are nonnormal nonisolated germs: the normalization fails along t=0, a locus of dimension d-1, because t remains integral but missing at its generic point. The result goes beyond normal smooth finite quotients; it is not inferred from characteristic-zero resolution of these singular germs.

For d=2 the source is a normal surface, so the same all-continuous-obstacle theorem holds over EVERY field, including imperfect fields and positive characteristic. There is also a mixed-characteristic example R=Z\_p[[t\textasciicircum{}2,t\textasciicircum{}3]]=Z\_p[[a,b]]/(b\textasciicircum{}2-a\textasciicircum{}3), S=Z\_p[[t]]. Grouping even/odd powers proves the indicated embedding and finiteness; S is complete regular of dimension two, its fraction field agrees, and t\textasciicircum{}2 S is the conductor. The reduced conductor quotients are both Z\_p. They have a nonempty centred boundary (the p-order ray), so infinite-conductor supports have not been discarded. N2 plus the normal-surface theorem proves CE for EVERY continuous signed obstacle here. Higher-dimensional positive/mixed-characteristic smooth source CE remains unproved by this argument.

For H=\textless{}2,3\textgreater{} there is a new explicit all-support example in EVERY dimension d\textgreater{}=2 and EVERY field, independent of any general source CE theorem. Use gauge a=(t\textasciicircum{}2,t\textasciicircum{}3,z\_1,...,z\_(d-1)), b=z\_1+t and J=(b,t\textasciicircum{}2,z\_2,...,z\_(d-1)) in S. J is primary. By the homeomorphism define the continuous downstairs A(v)=w(J), where rho(w)=v, and obstacle f=-A. For integers m\textgreater{}=1 form the ACTUAL downstairs primary ideal
K\_m=(b\textasciicircum{}m(z\_1-mt), t\textasciicircum{}2 b\textasciicircum{}m, t\textasciicircum{}(2m+2))+(z\_2,...,z\_(d-1))\textasciicircum{}m.
When d=2 omit the last summand and put U=+infinity below. The first generator has no t-linear coefficient: m z\_1\textasciicircum{}m-m z\_1\textasciicircum{}m=0, including in characteristic dividing m. All its other t-degrees are zero or \textgreater{}=2, so it lies in R. The remaining generators lie there too. The radical contains t\textasciicircum{}2,t\textasciicircum{}3,z\_1 and all the other z\_j, hence K\_m is proper primary. No generator is treated as a supplied executable expression.

Here is the universal bound, including EVERY support and infinite coordinate. Write s=w(t), u=w(z\_1), B=w(b), U=min\_(j\textgreater{}=2)w(z\_j), so min(2s,u,U)=1 and A=min(B,2s,U). A is finite. Every displayed generator has value \textgreater{}=mA, hence w(K\_m)\textgreater{}=mA. If U=A, the added ideal has value mU=mA, giving equality. Otherwise A\textless{}U. If 2s\textless{}=B, then s is finite and u=s: unequal values would force B=min(s,u)\textless{}2s. Since U\textgreater{}2s and min(2s,s,U)=1, s=u=1. The last original generator then has value mA+2. In the remaining case B\textless{}2s, A=B\textless{}U. If s\textless{}=u, then s is finite; B\textgreater{}=s. If s\textgreater{}1 all of 2s,u,U would exceed one, a contradiction. Thus s\textless{}=1, and t\textasciicircum{}2 b\textasciicircum{}m has value mA+2s\textless{}=mA+2. If s\textgreater{}u (including s=infinity), B=w(z\_1-mt)=u: for m nonzero in k, m is a unit and the values are unequal, and if m=0 in k the second expression is simply z\_1. The normalization and U\textgreater{}u give u=1, so the first generator has value mA+1. These cases exhaust all possibilities, including B=infinity, U=infinity, killed coordinates and cancellation. Consequently
0\textless{}=v(K\_m)/m-A(v)\textless{}=2/m on the ENTIRE normalized boundary.

Therefore f is a uniform limit from below of -v(K\_m)/m. Each profile is the evaluation of ACTUAL -E\_(K\_m)/m on ACTUAL Bl\_(K\_m)(R), negative effective and relatively ample, with positive degree on EVERY contracted integral curve. Thus f belongs to independently geometric C even in higher positive characteristic where general CE is unknown. It is its own greatest C/D minorant. Since K\_m subset A\_m\textasciicircum{}R(f), the canonical profiles themselves obey the EVERY-degree uniform estimate 0\textless{}=p(f)-h\_m\textasciicircum{}R(f)\textless{}=2/m. This establishes a quantitative geometric envelope across the infinite-conductor stratum without a conductor-value/N error or a root-map assumption.

In characteristic ZERO f is outside the ENTIRE finite REAL raw negative ideal cone Q\_R, and its canonical algebra is NOT finitely generated. Define w\_infinity by killing t,z\_2,... and giving z\_1 value one; define w\_star by killing z\_2,... and using monomial weights (t,b)=(1,2), with z\_1=b-t. Both satisfy the chosen gauge, and A(w\_infinity)=1, A(w\_star)=2. For ANY proper primary downstairs ideal I, its image after killing the other z\_j is primary in k[t\textasciicircum{}2,t\textasciicircum{}3,z\_1]. Let e=w\_infinity(I)\textgreater{}0. A generator attaining e has a unit denominator at both valuations. Its formal expansion has no t-degree-one terms, has nonzero z\_1\textasciicircum{}e coefficient, and has zero t z\_1\textasciicircum{}(e-1) coefficient. Its total-degree-e homogeneous part therefore cannot be a scalar multiple of b\textasciicircum{}e in characteristic zero. After z\_1=b-t it contains some t\textasciicircum{}(e-j)b\textasciicircum{}j with j\textless{}e, of w\_star-value e+j\textless{}=2e-1. Other total degrees cannot cancel this monomial. Thus w\_star(I)\textless{}2w\_infinity(I). Every nonzero finite nonnegative REAL combination of proper primary gauges preserves this strict inequality; unit terms vanish. A has equality, so A is outside the raw positive cone and f outside Q\_R. Finite homogeneous generation of the canonical algebra would force a rational primary envelope by the elementary argument in n1, contradicting this separation. No claim of non-finite-generation in positive characteristic follows from the characteristic-zero argument.

This is continuous-envelope regularity WITH a canonical O(1/m) bound and WITHOUT canonical finite generation for an explicit cancellation-sensitive nonnormal higher-dimensional example. It is not a counterexample to CE, nor a raw geometric-core attainment claim. Under another gauge the same homogeneously transported obstacle keeps all canonical ideals, and the bound is divided by min v(new gauge)\textgreater{}0. G331's literal radial scope distinction, curve min(f,0), nilpotent reduction and empty dimension-zero case remain unchanged.

Own finite\_diagnostics.py samples only its explicitly recorded rational weights, coherent cancellation cases and infinite-coordinate endpoints, and checks the t-linear cancellation exactly. Its output is finite diagnostic evidence, not the universal proof above. Artifacts n3.md, n3.tex, diagnostics-input.json and reports/finite-diagnostics.json; rerun python finite\_diagnostics.py; python assemble.py; python verify.py.

\subsection*{Expanded higher-dimensional cusp construction}
For a numerical semigroup $H$ of conductor $c$,
\[
R_H=k[t^H,z_1,\ldots,z_{d-1}]_0,\quad
S=k[t,z_1,\ldots,z_{d-1}]_0,\quad\mathfrak c=t^cS,
\]
\[
(R_H/\mathfrak c)_{\rm red}=(S/\mathfrak c)_{\rm red}=k[z_1,\ldots,z_{d-1}]_0.
\]
For $H=\langle2,3\rangle$, $b=z_1+t$ and
$A=w((b,t^2,z_2,\ldots,z_{d-1}))$, put
\[
K_m=(b^m(z_1-mt),\ t^2b^m,\ t^{2m+2})+(z_2,\ldots,z_{d-1})^m.
\]
The proof covers all supports and establishes
\[
0\leq \frac{v(K_m)}m-A\leq\frac2m,\qquad
0\leq p(-A)-h_m^R(-A)\leq\frac2m.
\]
In characteristic zero the two test points give $A(w_\infty)=1$,
$A(w_*)=2$, whereas every nonzero finite real positive primary combination
satisfies $u(w_*)<2u(w_\infty)$. Hence $-A\notin Q_R$ and the canonical
algebra is not finitely generated. Neither conclusion is a CE failure.

